% Русская рукопись. Сборка: Tectonic или XeLaTeX + BibTeX + XeLaTeX x2.
\documentclass{article}
\PassOptionsToPackage{numbers,sort,compress}{natbib}
\usepackage[preprint]{neurips_2025}
\usepackage{fontspec}
\setmainfont{cmunrm.otf}[BoldFont=cmunbx.otf,ItalicFont=cmunti.otf,BoldItalicFont=cmunbi.otf]
\setsansfont{cmunss.otf}[BoldFont=cmunsx.otf,ItalicFont=cmunsi.otf,BoldItalicFont=cmunso.otf]
\usepackage{polyglossia}
\setdefaultlanguage{russian}
\setotherlanguage{english}
\usepackage{amsmath,amssymb,amsfonts}
\usepackage{amsthm}
\usepackage{booktabs,tabularx,array}
\usepackage{graphicx}
\usepackage{microtype}
\usepackage{xcolor}
\usepackage{url}
\usepackage{hyperref}
\hypersetup{unicode=true,colorlinks=true,linkcolor=blue,citecolor=blue,urlcolor=blue}
\urlstyle{same}
\makeatletter
\renewcommand{\@noticestring}{}
\makeatother
\renewenvironment{abstract}{\vskip0.075in\centerline{\large\bfseries Аннотация}\vspace{0.5ex}\begin{quote}}{\par\end{quote}\vskip1ex}
\raggedbottom
\newcommand{\Pari}{P_{\mathrm{ari\_max}}}
\newcommand{\Psilh}{P_{\mathrm{silh\_max}}}
\newcommand{\ARI}{\operatorname{ARI}}
\newcommand{\sd}{\operatorname{sd}}
\newcolumntype{Y}{>{\raggedright\arraybackslash}X}
\title{Выбор кластерного разбиения по ограниченной внешней информации}
\author{Богданов А. Д.\\МГУ им. Ломоносова \And Сенько О. В. \\МГУ им. Ломоносова}
\begin{document}
\maketitle
\begin{abstract}
Работа посвящена выбору результата кластеризации при доступности меток лишь небольшой части объектов. В этой ситуации внешний критерий чувствителен к составу размеченной выборки, а внутренний индекс описывает геометрию данных, которая может не совпадать со структурой целевых классов. Исследуется, при каких условиях внутренний индекс улучшает выбор по ограниченной разметке и когда его использование приводит к потере качества. Рассматривается фиксированный пул разбиений и правило, которое допускает альтернативы лидеру по частичному скорректированному индексу Рэнда с учётом бутстрепной изменчивости парных различий, после чего выбирает кандидата по внутреннему индексу. Для оценки такого вмешательства задаётся протокол сравнения на скрытых проверочных метках с анализом внешнего качества, потерь относительно лучшего кандидата и последствий смены выбранного разбиения. Основные исследуемые факторы --- объём разметки, покрытие классов, согласованность геометрии с целевыми метками и состав пула кандидатов. Цель работы состоит в определении границ применимости совместного использования внутренних и внешних оценок, включая условия систематического ухудшения выбора.

\end{abstract}
\noindent\textbf{Ключевые слова:} кластеризация, внутренняя и внешняя валидация, ограниченная разметка, бутстреп, выбор разбиения, покрытие классов.
\section{Введение}\label{sec:intro}

Алгоритмы кластеризации с различными настройками могут формировать существенно разные группировки одних и тех же объектов. Поэтому прикладная задача включает не только построение разбиений, но и выбор результата, соответствующего цели анализа. Такой выбор особенно затруднён, когда известные метки описывают лишь небольшую часть данных: их недостаточно для надёжного сравнения всех кандидатов, однако они содержат информацию о том, какие различия между объектами значимы для задачи.

Для оценки разбиений используются два источника информации. Внутренние индексы характеризуют структуру групп в пространстве признаков. Например, silhouette сопоставляет расстояния внутри кластера и до соседних кластеров~\citep{rousseeuw1987}. Внешние критерии сравнивают разбиение с целевыми метками. Скорректированный индекс Рэнда (Adjusted Rand Index, ARI) учитывает согласие пар объектов с поправкой на случайное совпадение~\citep{hubert1985}. Высокая внутренняя оценка не гарантирует высокого внешнего качества: компактные группы могут объединять разные классы, а один класс может занимать несколько разделённых областей. В работе метки задают прикладной ориентир.

При малой разметке возникает дополнительная неопределённость. Разбиение, максимизирующее ARI на известных метках, может уступать другому разбиению на остальных объектах. При этом важен не только бюджет разметки, но и её состав. Две выборки одинакового размера могут представлять разные классы и содержать разное количество объектов редких категорий. Под покрытием классов далее понимается доля целевых классов, представленных хотя бы одной известной меткой. Эта характеристика используется для последующего анализа эксперимента. Её истинное значение не предполагается доступным правилу выбора.

Ограниченная внешняя информация уже использовалась для выбора моделей полуобучаемой кластеризации и готовых разбиений~\citep{pourrajabi2014,cobs2017}. Исследовались также способы объединения внешних оценок со структурными характеристиками кластеров~\citep{gertrudes2019}. Поэтому исследовательский вопрос заключается в условиях полезности такого объединения. В частности, необходимо установить, помогает ли учёт изменчивости частичного ARI определить, когда внутреннему индексу следует разрешить изменить выбор внешнего критерия.

Рассматриваемое правило сначала выбирает лидера по известным меткам. Для каждого альтернативного разбиения оценивается бутстрепный разброс разности его внешней оценки и оценки лидера. Кандидаты с небольшим отставанием относительно этого разброса образуют допустимое множество. Окончательный выбор определяется внутренним индексом. Такое правило связывает степень вмешательства геометрии с изменчивостью внешнего сравнения. Однако большой разброс может допустить существенно худшие кандидаты, а малый разброс не исключает ошибки при непредставленных классах. Кроме того, лидер выбирается на тех же метках, что используются для оценивания его отрыва, поэтому бутстрепный допуск не интерпретируется как доверительное множество~\citep{efron2014}.

Цель работы --- определить условия, при которых внутренний индекс снижает или увеличивает потерю внешнего качества при выборе из фиксированного пула. Исследование включает три задачи. Во-первых, сопоставить адаптивный допуск с выбором по отдельным критериям и простыми правилами их сочетания. Во-вторых, раздельно оценить влияние объёма разметки, покрытия классов и состава кандидатов. В-третьих, разобрать случаи смены внешнего лидера: как часто происходит вмешательство, какие альтернативы допускаются и к каким последствиям оно приводит на скрытых метках. Предполагаемый вклад состоит в описании границ применимости этого подхода и причин его ошибок. Возможное превосходство конкретного правила проверяется экспериментально.

\section{Связанные исследования}\label{sec:related}

\paragraph{Выбор разбиения по ограниченной внешней информации.}
В полуобучаемой кластеризации известные метки или парные ограничения могут использоваться как при построении кластеров, так и при выборе готового результата. Pourrajabi и соавторы~\citep{pourrajabi2014} рассматривают выбор модели и организацию проверки с учётом зависимости ограничений. COBS~\citep{cobs2017} использует ограничения для выбора между результатами разных алгоритмов и настроек, полученными без учителя. Авторы показывают экспериментально, что такой выбор способен превосходить сравниваемые полуобучаемые методы. Это непосредственный предшественник постановки с фиксированным пулом: внешняя информация направляет выбор, не изменяя самих кандидатов. В настоящей работе внешним сигналом служит ARI на подмножестве объектов, а предметом анализа --- последствия допуска альтернатив с учётом изменчивости этого сигнала.

\paragraph{Сочетание критериев качества.}
Akhanli и Hennig~\citep{akhanli2020} рассматривают внутренние индексы, характеризующие различные свойства кластеров, и калибруют их относительно случайных разбиений тех же данных перед агрегацией. Предложенные агрегированные индексы сравниваются на синтетических и реальных данных. Этот подход показывает, как сделать сопоставимыми оценки разных шкал, но не оценивает неопределённость внешнего критерия по малой разметке. Более близкий по источникам информации подход представлен в работе Castro Gertrudes и соавторов~\citep{gertrudes2019}: при извлечении кластеров из плотностной иерархии сочетаются структурная устойчивость и внешняя оценка на основе меток, включая неполное представление категорий. Следовательно, объединение внутренней и внешней информации и проблема отсутствующих классов имеют прямых предшественников. Здесь рассматривается выбор среди готовых разбиений разных алгоритмов и отдельно проверяется роль бутстрепного допуска.

\paragraph{Зависимость внутренних индексов от условий задачи.}
Benchmark-исследование Simpson и соавторов~\citep{simpson2026benchmark} сравнивает 26 внутренних индексов с использованием нескольких взаимодополняющих протоколов. Сопоставление выбранного лучшего разбиения и ранжирования кандидатов позволяет анализировать разные аспекты поведения индекса. В работе об instance space~\citep{simpson2026instance} девять индексов исследуются на 18\,351 синтетическом наборе и результатах восьми алгоритмов. Выявленные области успешной работы индексов мотивируют изучение условий их применимости вместо единого общего рейтинга. Для нашей задачи дополнительно существенны доступность внешних меток и величина потери после окончательного выбора: сходство ранжирований само по себе не определяет последствия конкретного решения.

\paragraph{Неопределённость после выбора.}
Efron~\citep{efron2014} рассматривает оценивание точности после выбора модели. Эта работа служит методологическим основанием учитывать влияние отбора на статистический вывод, но не обосновывает используемый здесь порог. Бутстреп относительно лидера, выбранного на той же размеченной выборке, характеризует условную изменчивость сравнения. Он не гарантирует правильности выбора при смещённой разметке и не восстанавливает отсутствующие классы.

Перечисленные направления определяют контекст исследования: выбор по ограниченной внешней информации, сочетание разных характеристик качества и анализ условий полезности индексов уже изучаются. Настоящая работа сосредоточена на проверке адаптивного допуска и диагностике его вмешательств в выбор внешнего лидера. Обзор имеет тематический характер и не является доказательством отсутствия аналогичного правила во всей литературе.

\section{Постановка задачи}\label{sec:problem}

\subsection{Данные, кандидаты и доступная информация}
Пусть $X=(x_1,\ldots,x_n)$, $x_i\in\mathbb R^d$, --- фиксированный набор объектов, а $y=(y_1,\ldots,y_n)$ --- скрытые метки целевых классов. До раскрытия меток строится конечный пул $\mathcal P=\{P_1,\ldots,P_M\}$. Элемент пула представляет полное разбиение набора. Кандидат $P$ задаёт принадлежность к кластеру для каждого объекта. Номера кластеров номинальны. Пул может включать результаты разных алгоритмов, чисел кластеров и настроек. Одинаковые разбиения объединяются.

Множество индексов объектов разделяется на резерв $S$ и проверочную часть $T$, $S\cap T=\varnothing$. В одном повторении из $S$ раскрывается подмножество $L$, $|L|=m$. Правило выбора получает признаки всех объектов, пул и метки $y_L$, но не получает $y_T$ или метки нераскрытой части резерва. Предобработка признаков без меток и построение кандидатов на всём $X$ разрешены трансдуктивной постановкой. Известные связанные объекты объединяются в группы и при разделении перемещаются совместно.

\subsection{Внутренняя и внешняя оценки}
Основной внутренний индекс $I(P)$ --- средний silhouette~\citep{rousseeuw1987}. Для точки $i$ обозначим через $a_i$ среднюю дистанцию до остальных объектов своего кластера, а через $b_i$ --- минимальную среднюю дистанцию до другого кластера. Тогда
\begin{equation}
 s_i=\frac{b_i-a_i}{\max(a_i,b_i)},\qquad I(P)=\frac1n\sum_{i=1}^n s_i.
\end{equation}
Для одноэлементного кластера и случая $a_i=b_i=0$ используется соглашение $s_i=0$. Дополнительная геометрическая оценка --- индекс Davies--Bouldin~\citep{davies1979}. При выборе направление его оптимизации меняется на противоположное, поскольку меньшее значение предпочтительно.

Частичная внешняя оценка и проверочное качество определяются на разных объектах:
\begin{equation}\label{eq:two-ari}
 A_L(P)=\ARI(y_L,P_L),\qquad Q_T(P)=\ARI(y_T,P_T).
\end{equation}
Пусть $n_{ab}$ --- число размеченных объектов класса $a$, попавших в кластер $b$. $n_{a\cdot}$ и $n_{\cdot b}$ --- суммы по строкам и столбцам таблицы сопряжённости. Для $m\ge2$ положим $U=\sum_{a,b}\binom{n_{ab}}2$, $V=\sum_a\binom{n_{a\cdot}}2$, $W=\sum_b\binom{n_{\cdot b}}2$, $N=\binom m2$. ARI имеет вид~\citep{hubert1985}
\begin{equation}
 \ARI=\frac{U-VW/N}{(V+W)/2-VW/N}.
\end{equation}
В вырожденных случаях используются соглашения реализации scikit-learn, в том числе значение 1 для совпадающих тривиальных разбиений. Такие случаи отдельно диагностируются.

Если алгоритм оставляет шумовые объекты без назначения, для внешних индексов каждый такой объект рассматривается как отдельная одноэлементная группа. В реализованном протоколе внутренние показатели сначала рассчитываются на назначенных объектах. Для доли назначенных $q$ приняты преобразования $S\mapsto qS$ и $DB\mapsto DB/q$. Кандидаты с $q<0{,}7$ или менее чем двумя кластерами исключаются. Эти соглашения фиксируются до сравнения методов. Преобразование $qS$ при отрицательном silhouette увеличивает значение индекса при уменьшении $q$, поэтому его нельзя интерпретировать как универсальный штраф за шум. Влияние обработки шумовых объектов требует отдельной проверки чувствительности.

\subsection{Правило выбора с адаптивным допуском}
Внешний лидер определяется только по доступной разметке:
\begin{equation}\label{eq:leader}
 \Pari=\arg\max_{P\in\mathcal P}A_L(P).
\end{equation}
Равенства разрешаются фиксированным порядком канонических идентификаторов кандидатов. Для каждого $P$ вычисляется наблюдаемое отставание
\begin{equation}
 d_P=A_L(\Pari)-A_L(P).
\end{equation}
Из $L$ формируются $B$ бутстрепных выборок $L^{(1)},\ldots,L^{(B)}$: объекты выбираются с возвращением, а при наличии связанных объектов перевыбираются целые группы. Повторения объектов учитываются с их кратностями. Пары объектов независимо не перевыбираются. Кластеризации и внутренние индексы остаются фиксированными.

Для всех кандидатов используются одни и те же бутстрепные выборки. Лидер из~\eqref{eq:leader} фиксируется на исходном $L$ и не переопределяется при вычислении
\begin{equation}
 d_P^{(b)}=A_{L^{(b)}}(\Pari)-A_{L^{(b)}}(P),\qquad
 \widehat\sigma_P=\sd_{b=1}^B(d_P^{(b)}).
\end{equation}
При $B\ge2$ оценка разброса задаётся явно:
\begin{equation}
 \bar d_P=\frac1B\sum_{b=1}^B d_P^{(b)},\qquad
 \widehat\sigma_P=\sqrt{\frac1{B-1}\sum_{b=1}^B(d_P^{(b)}-\bar d_P)^2}.
\end{equation}
Одна перевыборка используется для обоих разбиений в разности, чтобы сохранить зависимость их оценок на одних объектах. При $c\ge0$ задаются допустимое множество и выбранное разбиение:
\begin{equation}\label{eq:adaptive}
 \mathcal E_c=\{P\in\mathcal P:d_P\le c\widehat\sigma_P+\varepsilon\},\qquad
 \Psilh=\arg\max_{P\in\mathcal E_c}I(P),
\end{equation}
где $\varepsilon=10^{-12}$ --- численный допуск. Основное значение $c=1$.

Индекс $\mathrm{silh\_max}$ обозначает максимум silhouette внутри допустимого множества $\mathcal E_c$. Из определения следует $A_L(\Psilh)\le A_L(\Pari)$, однако порядок $Q_T(\Psilh)$ и $Q_T(\Pari)$ может быть любым. Этим объясняется возможность положительного эффекта на скрытой проверочной части без противоречия максимальности~\eqref{eq:leader}.

Правило~\eqref{eq:adaptive} не является тестом равенства качества. Бутстреп условен относительно наблюдённой разметки, а отбор лидера использует ту же выборку. Для множества $\mathcal E_c$ не заявляются доверительное покрытие или известная вероятность ошибки~\citep{efron2014}. При отсутствии классов в $L$ их информация не появляется ни в одной перевыборке.

\subsection{Покрытие и оцениваемый эффект}
Истинное покрытие классов определяется ретроспективно:
\begin{equation}
 \operatorname{cov}(L)=
 \frac{|\{y_i:i\in L\}|}{|\{y_i:i=1,\ldots,n\}|}.
\end{equation}
Оно используется для объяснения поведения методов, но не для вычисления~\eqref{eq:adaptive}. Покрытие классов отличается от доли назначенных объектов $q$ в плотностной кластеризации.

После фиксации решений раскрываются проверочные метки. Основной показатель --- потеря относительно лучшего кандидата того же пула:
\begin{equation}
 R_T(\widehat P)=\max_{P\in\mathcal P}Q_T(P)-Q_T(\widehat P)\ge0.
\end{equation}
Отдельно учитывается абсолютное качество: малая потеря в слабом пуле не означает хорошей кластеризации. Эффект адаптивного выбора относительно внешнего лидера равен
\begin{equation}
 G_T=Q_T(\Psilh)-Q_T(\Pari)
     =R_T(\Pari)-R_T(\Psilh).
\end{equation}
Положительное значение означает улучшение. Для диагностики вводится индикатор переключения
\begin{equation}
 Z_L=\mathbf 1\{\Psilh\ne\Pari\}.
\end{equation}
Если $Z_L=0$, то $G_T=0$. Обратное неверно: разные разбиения могут иметь одинаковый проверочный ARI. Поэтому частота переключений оценивается отдельно от долей улучшений и ухудшений. При положительной вероятности переключения
\begin{equation}
 \mathbb E[G_T]=\Pr(Z_L=1)\,\mathbb E[G_T\mid Z_L=1],
\end{equation}
где ожидание относится к зафиксированной схеме раскрытия меток при заданных данных, пуле и проверочной части. Эта запись разделяет частоту вмешательства и его среднее последствие. Для фиксированного набора исследуются среднее $G_T$ по повторениям разметки и доли улучшений, совпадений и ухудшений. Эти повторы не являются независимыми наборами данных.

Базовые сравнения включают максимум внутреннего индекса без меток, внешний максимум, геометрическое разрешение равенств, фиксированную смесь, смесь с весом, замороженным на отдельной development-серии, фиксированный допуск по частичному ARI и выбор по частоте бутстрепных побед без внутреннего индекса. В коде им соответствуют M0--M7. $\Pari$ и $\Psilh$ соответствуют M1 и M6. Настройка по проверочным меткам подтверждающей серии не допускается.

\subsection{Экспериментальные условия и проверяемые гипотезы}
Основные факторы --- бюджет $m$, схема раскрытия меток и структура данных. Случайный сценарий выбирает объекты равномерно без возвращения. В сценарии отсутствующих классов метки раскрываются только из случайно выбранной половины категорий, хотя все объекты остаются в $X$. В смещённом сценарии половина классов получает больший вес отбора. Бюджеты внутри повторения вложены. Все методы получают одно и то же $L$.

Первый этап исследования --- диагностический анализ сохранённых пилотов. Для каждого повторения сопоставляются внешний лидер и окончательный выбор, их частичный и проверочный ARI, внутренние оценки, размер допустимого множества, отставание и бутстрепный разброс. Сравнение с разрешением только равенств, фиксированным допуском и бутстрепным выбором без внутреннего индекса позволяет отделить эффект геометрии от эффекта самого допуска. Повторные вычисления при разных начальных состояниях генератора бутстрепа оценивают чувствительность решений к конечному $B$.

После диагностики предусматривается расширенная серия с независимыми синтетическими реализациями и несколькими реальными числовыми наборами. При одинаковом бюджете отдельно варьируются покрытие классов и распределение меток между ними. Состав пула и соответствие геометрии целевым классам рассматриваются как самостоятельные факторы. Параметры правила определяются на отдельной настроечной серии и фиксируются до анализа проверочной серии. Существующие пилоты, использованные для разработки правил, не считаются независимым подтверждением их эффективности.

Проверяются следующие предположения: (H1) польза внутреннего индекса зависит от соответствия геометрии целевым классам. (H2) влияние бюджета меток нельзя объяснить без учёта их состава. (H3) при неполном покрытии малый бутстрепный разброс может сопровождаться существенной ошибкой выбора. Основные результаты сравнения должны включать среднюю потерю, частоту и последствия переключений, вероятность ухудшения и величину отрицательных эффектов. Повторы разметки характеризуют вариативность внутри задачи. Переносимость выводов оценивается по независимым задачам. Ни одна из гипотез не предполагает универсального преимущества адаптивного правила.

\bibliographystyle{unsrtnat}
\bibliography{references}
\end{document}
